\documentclass[11pt,a4paper]{article}
\usepackage[margin=24mm]{geometry}
\usepackage[T1]{fontenc}
\usepackage{lmodern,microtype,amsmath,amssymb,amsthm,mathtools,mathrsfs}
\usepackage{booktabs,longtable,array,enumitem,xcolor,xurl,hyperref,fancyhdr}
\hypersetup{colorlinks=true,linkcolor=blue!45!black,citecolor=blue!45!black,urlcolor=blue!45!black,
 pdftitle={Congruence separation and positive-word fusion for explicit consistency equations},
 pdfauthor={AI-assisted research artifact prepared for ProveIt}}
\setlength{\emergencystretch}{2em}
\setlength{\parskip}{4pt}
\setlength{\parindent}{0pt}
\setlist{itemsep=2pt,topsep=4pt}
\allowdisplaybreaks
\setcounter{tocdepth}{1}
\pagestyle{fancy}\fancyhf{}\fancyhead[L]{\small Explicit consistency equations}
\fancyhead[R]{\small 23 October 2026}\fancyfoot[C]{\thepage}
\setlength{\headheight}{14pt}
\newtheorem{theorem}{Theorem}[section]
\newtheorem{lemma}[theorem]{Lemma}
\newtheorem{proposition}[theorem]{Proposition}
\newtheorem{corollary}[theorem]{Corollary}
\theoremstyle{definition}\newtheorem{definition}[theorem]{Definition}
\theoremstyle{remark}\newtheorem{remark}[theorem]{Remark}
\newcommand{\Z}{\mathbb Z}\newcommand{\N}{\mathbb N}\newcommand{\Pos}{\mathbb Z_{>0}}
\newcommand{\Sq}{\operatorname{Sq}(\Z)}\newcommand{\Con}{\operatorname{Con}}
\newcommand{\pop}{\operatorname{pop}}\newcommand{\lc}{\operatorname{lc}}
\newcommand{\code}[1]{\texttt{\small\detokenize{#1}}}
\newcommand{\TenDegree}{9,152,284}
\newcommand{\TenOps}{586}
\newcommand{\NineDegree}{5,671,669}
\newcommand{\NineOps}{630}
\newcommand{\NineNatOps}{636}
\newcommand{\FastDegree}{6,692,108}
\newcommand{\FastOps}{604}
\newcommand{\FastNatOps}{610}

\title{\textbf{Congruence separation and positive-word fusion}\\[3pt]
\large Smaller explicit Diophantine consistency equations for ProveIt}
\author{AI-assisted research artifact prepared for the ProveIt project\\
\small Mathematical constructions, literal circuits, and independent finite audits}
\date{23 October 2026\\\small Repository baseline: \texttt{4b4349064a75}}

\begin{document}
\maketitle
\begin{abstract}
We construct explicit integer polynomials whose unsolvability is equivalent
in ZFC to the syntactic consistency of ZFC, by improving the arithmetic
reductions attached to ProveIt's fixed refutation-search compiler. The main
construction has ten unrestricted integer variables, exact total degree
\TenDegree, and a division-free circuit of \TenOps\ binary arithmetic
operations. The strongest incoming construction at the pinned repository
snapshot has degree 21,111,400 and 659 operations in the same domain.
The principal new step combines two square conditions and a positivity
condition in one coprime-factor guard. Congruences distinguish that radicand
from the remaining Lucas radicand, so a four-factor norm suffices.
For nine positive variables we obtain degree \NineDegree\ with \NineOps\
operations; a cheaper variant has degree \FastDegree\ with \FastOps\
operations. A secondary fourteen-positive-variable construction has degree
52 and 79 operations. A common source transformation reduces 136 weighted
records to 94, with a newly checked selective code. The archive includes
all literal circuits, coefficient payloads and bounded recipes, proofs,
generators, and independent finite checks. The reductions are ordinary
mathematics formalizable in ZFC; no new end-to-end proof-assistant certificate
or global minimality claim is asserted.
\end{abstract}
\tableofcontents
\newpage

\section{Results, objective, and logical meaning}

There is no unique mathematical ordering called ``simplest Diophantine
equation.'' Number of supplied variables, their domain, total degree,
straight-line evaluation cost, coefficient height, and proof complexity are
different criteria. We retain this distinction and improve concrete points
of the supplied repository's frontier. We do not claim a lower bound or an
optimal equation among all integer polynomials.

\begin{theorem}[Main constructions]\label{thm:main}
Let $\Con(\mathrm{ZFC})$ mean that no finite proof in the repository's fixed
recursive presentation of ZFC derives a contradiction. The literal circuits
specified in this report and the accompanying data define fixed polynomials
$F_{10}\in\Z[\mathbf t]$, $F_9,F_{9,\mathrm f}\in\Z[\mathbf x]$, and
$P_{14}\in\Z[\mathbf u]$, with respectively ten, nine, nine, and fourteen
variables. Composing the proved arithmetic reductions below with the
inherited calibrated source theorem gives
\begin{align*}
\mathrm{ZFC}&\vdash
 \bigl(\forall\mathbf t\in\Z^{10}\;F_{10}(\mathbf t)\ne0\bigr)
 \ \longleftrightarrow\ \Con(\mathrm{ZFC}),\\
\mathrm{ZFC}&\vdash
 \bigl(\forall\mathbf x\in\Pos^9\;F_9(\mathbf x)\ne0\bigr)
 \ \longleftrightarrow\ \Con(\mathrm{ZFC}),
\end{align*}
and the same statement holds for $F_{9,\mathrm f}$ and for $P_{14}$ over
$\Pos^{14}$. Their exact degrees and emitted circuit costs are those in
Table~\ref{tab:results}.
\end{theorem}

The phrase ``inherited calibrated source theorem'' is substantive. It is
the fixed positive-source-to-blank-halting equivalence and the calibration
of the literal K215 machine to syntactic ZFC refutations, not an assumption
that ZFC is consistent. Section~\ref{sec:status} identifies the proof
dependencies and the separate formal-verification boundary. Each new
arithmetic reduction can also be read as an unconditional theorem about
existence of zeros of the explicitly retained positive source.

\begin{table}[ht]
\centering\small
\begin{tabular}{@{}lrrr@{}}
\toprule
Construction and supplied domain & Variables & Exact degree & Operations\\
\midrule
New two-root integer construction & 10 in $\Z$ & \TenDegree & \TenOps\\
Incoming selective integer baseline & 10 in $\Z$ & 21,111,400 & 659\\
\addlinespace
New majorant construction & 9 in $\Pos$ & \NineDegree & \NineOps\\
New cheaper majorant variant & 9 in $\Pos$ & \FastDegree & \FastOps\\
Incoming selective positive baseline & 9 in $\Pos$ & 7,344,328 & 738\\
\addlinespace
New auxiliary-gap construction & 14 in $\Pos$ & 52 & 79\\
Incoming support-sensitive baseline & 14 in $\Pos$ & 56 & 79\\
\bottomrule
\end{tabular}
\caption{Like-for-like comparisons at commit \texttt{4b4349064a75}.
Every operation count is for a particular emitted live circuit, not a
minimal-circuit theorem.}\label{tab:results}
\end{table}

The ten-integer degree decreases by approximately 56.65\%, with approximately
11.08\% fewer operations. The nine-positive majorant degree decreases by
approximately 22.77\%, with approximately 14.63\% fewer operations.
The number of variables is unchanged in each comparison. The nine-variable
results also have natural-coordinate versions: replacing each positive
coordinate by a natural coordinate plus one preserves the degree. The
emitted natural circuits use \NineNatOps\ and \FastNatOps\ operations.

These comparisons are against the actual incoming archives, including the
latest selective package, rather than only the older repository summary.
Other existing points remain useful: the incoming sparse-Pell package has
a nine-positive circuit of degree 15,502,748 and 548 operations, and a
ten-integer circuit of degree 50,822,936 and 582 operations. The incoming
short constructions include fifteen-positive circuits of degree 44 with
75 operations and degree 157 with 66 operations. Thus the present work does
not establish the fewest-operation record over all variable counts.

An operation is a live binary addition, subtraction, or multiplication;
squaring and multiplication by a numeral are charged. Shared gates are
charged once. Literal and fixed-coefficient leaves are not charged. Fixed
coefficient generation and the bit cost of enormous integers are reported
separately. No variable exponentiation or division occurs in any final
circuit. The full monomial expansions would be unhelpfully large; a finite
acyclic arithmetic circuit is an exact definition of an integer polynomial.

The theorem proves an equivalence, not either of its two sides. In
particular it is not a proof of $\Con(\mathrm{ZFC})$ or of unsolvability.
No consistency assumption is used to select the coefficient tuple. This
avoids both a nonconstructive ``choose zero or one depending on consistency''
definition and any conflict with the second incompleteness theorem.

\section{The fixed source and positive-word fusion}\label{sec:source}

\subsection{The retained positive quadratic system}

Let $p_1,\ldots,p_{22}$ be the original package-31 residuals in 31 positive
coordinates, listed in Appendix~\ref{app:source}. Their four fixed ports
are $(\beta,K_0,M_C,M_F)$, with $M_F$ already shifted as required by that
source. We use the whole \code{wang943-mixed-121x8801} tuple, never a mixture
of individual constants from different compilers. Its native radix is
$2^{1064921}$. The fixed integer bit lengths are
\[
(1064921,\ 1070004,\ 1064921,\ 1064922).
\]
The four unsigned little-endian payloads contain 533,099 bytes in total
before compression. The archive contains the payloads, their byte hashes,
their finite compiler recipe, and the literal machine and tile data from
which they are generated. In particular these are instantiated integers,
not four extra quantified variables.

The inherited source theorem is
\begin{equation}\label{eq:source-endpoint}
 \exists\mathbf z\in\Pos^{31}\;\bigwedge_{j=1}^{22}p_j(\mathbf z)=0
 \quad\longleftrightarrow\quad
 \exists n\;\operatorname{Refutation}_{\mathrm{ZFC}}(n).
\end{equation}
The unbounded native-window compiler and calibrated-machine proof are
retained dependencies. The new source arithmetic below proves equivalence
to the left-hand side of~\eqref{eq:source-endpoint}; it does not replace
that compiler with a finite simulation.

\subsection{A global shift of the packed word}

Write $p_j(\mathbf z)=\sum_{|\alpha|\le2}c_{j,\alpha}\mathbf z^\alpha$ and set
\[
 d_{j,\alpha}=(2-|\alpha|)!\prod_i\alpha_i!\,c_{j,\alpha}.
\]
Choose the weights $w_i$ and row targets $q_j$ of Appendix~\ref{app:code},
and define the fixed sparse polynomial
\begin{equation}\label{eq:Dcode}
 \mathcal D(T)=\sum_{j,\alpha}d_{j,\alpha}
 T^{q_j-\sum_iw_i\alpha_i},\qquad
 \ell(T)=\sum_{i=1}^{31}T^{w_i},\quad
 \mathcal A(T)=\sum_{j=1}^{22}T^{q_j}.
\end{equation}
All exponents in $\mathcal D$ are nonnegative. The largest coordinate
weight is $W=183$, the largest target is $Q_c=5190$, and
\[
 K_c=Q_c+2W=5556,\qquad E_c=K_c+W+2=5741.
\]
The certificate enumerates each of 94 supported source records against all
$\binom{33}{2}=528$ padded input pairs. Exactly 94 of the 49,632
interactions hit marked targets, and every hit belongs to its intended row
and monomial. Padding uses coordinate zero of weight zero. The factorial
factor above times the expansion multiplicity is always two. Consequently
\begin{equation}\label{eq:extraction}
 [T^{q_j}]\left(1+\sum_i z_iT^{w_i}\right)^2\mathcal D(T)
 =2p_j(\mathbf z).
\end{equation}
This is an exact polynomial identity, not an empirical statement about a
sample of large integers.

Previously each residual was expanded at $\mathbf n+\mathbf1$, producing
136 weighted records. We instead encode the nonnegative offsets by
$g=\sum_i n_iB^{w_i}$ and use
\begin{equation}\label{eq:word-shift}
 1+g+\ell(B)=1+\sum_i(n_i+1)B^{w_i}.
\end{equation}
The positive source can therefore remain unexpanded, with only 94 records.
This is a substitution in the common packing word rather than 22 separate
translations. There is an essential finite correction: translation had
canceled two monomials, so the old shifted-support code did not certify the
original support. We retarget row 13 from 2212 to 2124 and row 21 from 1050
to 1059. All other weights and targets remain unchanged. The legacy code
is retained as a negative control: it has five forbidden interactions on
the new support. The corrected code passes the complete enumeration.

\subsection{Bounds and the sparse borrow argument}

Put
\[
 M=2^{2134927},\qquad \kappa=2134940,\qquad a_c=2^\kappa=8192M,
 \qquad B=a_cb^2.
\]
The actual coefficient recipes give $|[T^r]\mathcal D(T)|<M$. Their largest
triangle bound has bit length 2,134,926. Independently checked support data
give $\deg\mathcal D=5102$ and leading coefficient one. Thus for $b\ge2$,
\[
 \mathcal D(B)>B^{5102}-M\frac{B^{5102}-1}{B-1}>0.
\]
The following common registers will be used by both few-variable routes:
\begin{align}
 N_0&=B^{184},& Z_c&=B^{5557},\notag\\
 M_c&=N_0-1-(b-1)\ell(B),&
 T_c&=M_c+N_0(B-2)\mathcal A(B),\label{eq:source-registers}\\
 J_c&=2(1+g+\ell(B))^2\mathcal D(B)+\mathcal A(B),&
 S_c&=g+N_0J_c.\notag
\end{align}
Under the larger of the two needed bounds, $0<g<4bB^{183}$,
the distinct positive weights imply $1+\ell(B)<bB^{183}$ and
$1+g+\ell(B)<5bB^{183}$. Hence
\[
 0<\frac{J_c}{Z_c}
 <\frac{50Mb^2+1}{B-1}<\frac1{128}.
\]
The last inequality reduces to $129<1792Mb^2$. Also
$0<T_c<N_0Z_c/128$ after the same elementary geometric estimates, and
$S_c<N_0Z_c/64$; the margins are more than sufficient for either block
length used below. These bounds precede digit decoding.

Suppose now that $b$ is a power of two. The low mask condition
$g\mathbin{\&}M_c=0$, with $g<N_0$, says exactly that
\[
 g=\sum_i n_iB^{w_i},\qquad 0\le n_i<b.
\]
The positive coordinates $z_i=n_i+1$ lie in $[1,b]$. Coefficients of
$(1+\sum z_iT^{w_i})^2\mathcal D(T)$ have absolute value less than
\[
 M(1+31b)^2<1024Mb^2=B/8.
\]
Raw coefficients of $J_c$ are therefore bounded in absolute value by
$B/2$. At target $q_j$ their exact value is $4p_j(\mathbf z)+1$.
Base-$B$ normalization has incoming carries only in $\{-1,0\}$: induction
uses $\lfloor(c+\varepsilon)/B\rfloor\in\{-1,0\}$ in this range.
The marked mask $B-2$ permits precisely digits zero or one. Such a digit
cannot have outgoing carry $-1$, since the required raw value would then
be near $-B$, outside the bound. The target digit is thus
$4p_j(\mathbf z)+1+\varepsilon$, with $\varepsilon\in\{-1,0\}$.
It is zero or one exactly when $p_j(\mathbf z)=0$. Positivity of $J_c$
ensures termination of the borrows. No infinite negative radix expansion
is being treated as a nonnegative word.

Conversely, given a positive source zero, choose an arbitrarily large
power of two $b$ and encode $n_i=z_i-1$. The word $g$ is positive:
the all-one source vector is not a zero, because residual 5 is one there,
independently of all fixed ports. It satisfies $g<bB^{183}$ and is divisible
by $B$, hence by $b^2$. These facts will justify the congruence restrictions
used later. They are properties of the constructive source witnesses, not
assumptions about arbitrary solutions of the final equations.

\section{Two algebraic tools}\label{sec:tools}

\begin{lemma}[Unweighted two-root norm]\label{lem:norm}
Let $a,d\in\Z$ with $a\ne d$. If $t\in\mathbb Q$ satisfies
\[
 (t^2-a-d)^2-4ad=0,
\]
then $a,d$ are integer squares and $t$ is an integer. Conversely any signed
integer roots $u^2=a,v^2=d$ give a zero with $t=u+v$.
\end{lemma}
\begin{proof}
At $t=0$ the left side is $(a-d)^2$, so $t\ne0$. Define
\[
 u=\frac{t^2+a-d}{2t},\qquad v=t-u.
\]
Direct expansion gives $u^2=a$ and $v^2=d$. A rational number whose square
is an integer is an integer, by writing it in lowest terms. Thus
$u,v,t$ are integers. The converse is the same identity.
\end{proof}

There is no premise about the sign of the candidate radicands and no
square-class independence premise. With exactly two radicals, an
unconditional inequality $a\ne d$ is enough. Replacing large numerical
separators by cheap congruences is the central design change.

\begin{lemma}[Coprime square and positivity guard]\label{lem:guard}
Suppose $b,Q\in\Z$, $|b|\ge2$, $\gcd(b,Q)=1$, and
$b\equiv Q\equiv1\pmod3$. Suppose also $Q\ne0$, and set $r=bQ$.
For an integer $O$ define
\[
 A_O(z)=r+(3O-1)(rz)^2.
\]
There is a nonzero integer $z$ making $A_O(z)$ a square if and only if
$b,Q$ are positive squares and $O>0$.
\end{lemma}
\begin{proof}
Factor $A_O(z)=r[1+(3O-1)rz^2]$. The factors are coprime and nonzero:
the second is one modulo $r$, and $|r|\ge2$. If their product is a square,
they are signed squares with the same sign. Since $r\equiv1\pmod3$, it
cannot be a negative square. Both factors are positive squares. Applying
the same coprime-factor argument to $r=bQ$, and again using residue one
modulo three, shows that $b,Q$ are positive squares. If $O\le0$, the
second factor is at most $1-r<0$, a contradiction.

Conversely, $r$ is a positive square, and $3O-1$ is a positive nonsquare,
being two modulo three. Thus $(3O-1)r$ is a positive nonsquare. Pell's
theorem supplies $u,z>0$ with $u^2-(3O-1)rz^2=1$. Then
$A_O(z)=ru^2$ is a square.
\end{proof}

\begin{lemma}[The retained quintic merger]\label{lem:quintic}
The polynomial $\pi(T)=(T^3-T-1)(T^2+23)$ has no rational root and has a
root modulo every positive integer. For $S\ne0$, put
\[
 \Theta=(\rho^3-\rho S^2-S^3)(\rho^2+23S^2),\qquad \rho=T+Sy.
\]
Then $\Theta\ne0$. If $\xi S^5\mid\Theta$, then $\xi\ne0$ and $S\mid T$.
When $S\mid T$ and $\xi\ne0$, the divisibility can be arranged by a
prescribed residue class of $y$ modulo $\xi$.
\end{lemma}
\begin{proof}
The rational-root theorem excludes a root of the cubic; the quadratic is
positive over $\mathbb Q$. For an odd prime $p\ne23$, either the cubic has
a simple root or it is irreducible. In the second case its Frobenius
permutation is a three-cycle, so its discriminant $-23$ is a square modulo
$p$; the quadratic then has a simple root. At $p=23$ the cubic has the
simple root 3. Hensel lifting handles odd prime powers. For powers of two,
$-23\equiv1\pmod8$ is an odd square in $\mathbb Z_2$; the elementary
successive lifting of odd squares gives the same conclusion. The Chinese
remainder theorem combines the roots.

Now $\Theta=S^5\pi(\rho/S)$ cannot vanish. If a prime satisfies
$\beta=v_p(T)<v_p(S)=\alpha$, then $v_p(\rho)=\beta$. The two factors of
$\Theta$ have valuations exactly $3\beta$ and $2\beta$, so
$v_p(\Theta)=5\beta<5\alpha$, contradicting $S^5\mid\Theta$.
Finally, when $T=Sq$, choose $c$ with $\xi\mid\pi(c)$ and impose
$y\equiv c-q\pmod\xi$.
\end{proof}

The quintic merger was already available in the supplied constructions;
it is retained here, not claimed as a new invention.

\section{The ten-integer equation}\label{sec:ten}

\subsection{All polynomial definitions}

The ten supplied variables are
\[
 (t,g_0,h_0,k_0,l,w,x,y,z,m)\in\Z^{10}.
\]
The circuit names the first coordinate \code{f}. Define
\begin{equation}\label{eq:normalization}
 b=6t+4,\quad g=bg_0,\quad h=bh_0,\quad k=b(3k_0+2),\quad B=a_cb^2.
\end{equation}
In the inherited bridge the old first coordinate is $f=2t+1$, so
$b=1+3f$. The emitted circuit evaluates $6t+4$ directly.
Use $S_c,T_c$ from~\eqref{eq:source-registers}, and set
\begin{align}
 N_c&=B^{5741},&
 X&=N_c^2+(S_c+T_c)N_c+N_c-1-S_c,\notag\\
 Y&=a_c^{15}\,2^{2134919}\,B^{5931}b,&
 \mathscr C&=bB^{183}.\label{eq:ten-source}
\end{align}
The following registers use the modified Lucas bridge of
Sun~\cite{sun2021,sun2026}:
\begin{align}
 L&=lY,& U&=2LX,& V&=4gwY,& W_*&=bw,\notag\\
 A&=U(V+1),& B_*&=2X+1,& C&=B_*+(A-2)h,\notag\\
 P&=U^2V,& K&=X+1+k(P-2),& S&=2A-5,\notag\\
 T&=3W_*C-2(W_*^2-1),&&
 Q=\bigl(P_{1/2}^2-1\bigr)K^2+1,\label{eq:ten-bridge}
\end{align}
where $P_{1/2}=(LX)UV$ is a polynomial and $2P_{1/2}=P$. In particular
$4Q=(P^2-4)K^2+4$, the bridge's auxiliary square expression. Put
\begin{equation}\label{eq:ten-margin}
 \Delta=8\mathscr C^3gK^2
 -g^2\left(32(C-KL)^2\mathscr C^3+g^2K^2\right).
\end{equation}
For $\xi=zx$ define
\begin{align}
 D&=(A^2-4)C^2+4,& E&=C^2D\xi,\notag\\
 F&=4(A^2-4)E^2+1,& G_*&=2(A+2)(A-2)^2E^2-1,\notag\\
 H_*&=C+(B_*-C)F+16EFy,& I_*&=(G_*^2-1)H_*^2+1.
 \label{eq:ten-lucas}
\end{align}
Finally let
\begin{align}
 \rho&=T+Sy,& \Sigma&=\xi S^5,&
 \Theta&=(\rho^3-\rho S^2-S^3)(\rho^2+23S^2),\notag\\
 r&=bQ,& a&=r+(3\Delta-1)(rz)^2,& d&=DFI_*,\notag\\
 Z&=\Sigma m-\Theta.&&&&\label{eq:ten-final-registers}
\end{align}
The requested single equation is
\begin{equation}\boxed{
 F_{10}=\left(Z^2-\Sigma^2(a+d)\right)^2-4\Sigma^4ad=0.
 }\label{eq:ten-final}
\end{equation}
Every letter on the right is a computed abbreviation. Only the ten listed
coordinates are unknowns. The equations defining the abbreviations do not
add existential witnesses.

\subsection{The exact imported interface}

For clarity, the number-theoretic input used here is stronger than a bare
existence theorem for some ten-variable representation. Sun's 2021
Theorems 4.12--4.13, specialized at $p=P_{\rm thm}=2$ and $Q_{\rm thm}=1$,
and the modified certificate and interpolation of the 2026 preprint give
the following interface.

From a valid source encoding with $b$ a power of two,
$0<g<\mathscr C$, and $Y\mid\binom{2X}{X}$, the constructive proof chooses
positive $l,w$ and integral $h,k$ such that
\[
 C=u_{2X+1}(A),\quad K=u_{X+1}(P),\quad
 Q\in\Sq,\quad S\mid T,\quad\Delta>0.
\]
Here $u_0(A)=0,u_1(A)=1,u_{n+1}(A)=Au_n(A)-u_{n-1}(A)$.
The exact $C$ and $K$ choices occur immediately before and in equation
(4.16) of the 2021 proof. They matter for the new residue restrictions.
The interpolation theorem allows $\xi$ to be divisible by any prescribed
positive integer and subsequently allows any residue class of $y$ modulo
$\xi$, while preserving $DFI_*\in\Sq$.

Conversely, for the displayed registers the hypotheses
\[
 b\in\Sq,\quad f^2l^2\Delta>0,\quad\xi\ne0,\quad
 4Q\in\Sq,\quad DFI_*\in\Sq,\quad S\mid T
\]
recover $b$ as a power of two, $0<g<2\mathscr C$, and the binomial
divisibility. This is the interface of the retained reverse bridge; its
source size conditions hold after the positive-word change.

Indeed, for $b=2^s$ and $B=2^v$, the explicit mask has
\[
 \pop(T_c)=206v-31s-22.
\]
The two-block definition of $X$, with $0\le S_c,T_c$ and $S_c+T_c<N_c$,
gives
\[
 \pop(X)=1+5741v+\pop(T_c)
 -\{\pop(S_c)+\pop(T_c)-\pop(S_c+T_c)\}.
\]
Since $v_2\binom{2X}{X}=\pop(X)$ and
$\log_2Y=5947v-31s-21$, divisibility by $Y$ is exactly absence of the
indicated binary carries. Section~\ref{sec:source} then decodes the source
zero. For the forward implication we may take any sufficiently large
$b=4^s$, so it is also a square. This explains why the more specialized
source still meets the imported interface.

\subsection{Universal congruences and radicand separation}

The following statements hold on every signed assignment, before
assuming the final equation or any positivity. We have $b\ne0$,
$|b|\ge2$, $2\mid b$, $b\equiv1\pmod3$, and $a_c\equiv1\pmod3$.
Both $N_0$ and $N_c$ are divisible by $b^2$, so
\[
 X+1\equiv-g\equiv0\pmod b.
\]
Also $X$ is odd. Modulo three, all powers of $B$ are one, $M_c=0$, and
$T_c\equiv-22\equiv2$. The terms involving $S_c$ cancel, giving
$X\equiv1+T_c\equiv0\pmod3$.

The factor $b$ in $Y$ gives $2b\mid U$ and $2b\mid A$. Hence $b\mid P$,
and $3\mid P$ follows from $3\mid X$. The imposed $k$ gives
\[
 b\mid K,\qquad3\mid K,\qquad
 Q\equiv1\pmod{b^2},\qquad Q\equiv1\pmod3.
\]
Thus $Q\ne0$, $Q$ is odd, and $\gcd(b,Q)=1$. Also
\[
 C+1=2(X+1)+(A-2)bh_0\equiv0\pmod{2b}.
\]
Writing $D=A^2C^2-4(C^2-1)$ shows $4b\mid D$. Therefore
\begin{equation}\label{eq:ten-distinct}
 \frac ab=Q+(3\Delta-1)bQ^2z^2\equiv1\pmod2,
 \qquad \frac db\equiv0\pmod2.
\end{equation}
The quotients in this proof are integers; they are not circuit operations.
As $b\ne0$, equation~\eqref{eq:ten-distinct} proves $a\ne d$
unconditionally. A large separating weight is unnecessary.

\subsection{Soundness, including all degenerate cases}

Suppose $F_{10}=0$. The integer $S=2A-5$ is odd and nonzero.
Lemma~\ref{lem:quintic} gives $\Theta\ne0$. If $\Sigma=0$, the final
polynomial would equal $\Theta^4$, a contradiction. Thus
$\Sigma\ne0$, in particular $\xi\ne0$ and $z\ne0$.

Divide the norm identity by $\Sigma^4$ and apply
Lemma~\ref{lem:norm} with $t=Z/\Sigma$. It yields $a,d\in\Sq$ and
$Z/\Sigma\in\Z$, hence $\Sigma\mid\Theta$. The quintic merger gives
$S\mid T$. Lemma~\ref{lem:guard} now gives $b,Q$ positive squares and
$\Delta>0$.

As $\mathscr C>0$, formula~\eqref{eq:ten-margin} forces $g>0$ and $K\ne0$:
if $g\le0$ every contribution is nonpositive, and if $K=0$ the remaining
term is nonpositive. If $l=0$, then $U=P=P_{1/2}=0$ and
$Q=1-K^2\le0$, again impossible. Thus $l\ne0$.
The old $f=2t+1$ never vanishes, so $f^2l^2\Delta>0$.
This recovers the inherited nonvanishing hypothesis although its
$f^2l^2$ factor has been omitted from the new guard. Since $Q$ is a square,
$4Q$ is a square. All hypotheses of the reverse interface now hold,
and it decodes a positive source zero. No prior sign assumption on $d$ or
on the free certificate variables was used.

\subsection{Completeness of the normalization}

Suppose the positive source has a zero. Choose sufficiently large
$b=4^s$ with $s\ge1$. Then $t=(b-4)/6$ is an integer, and its canonical
packed $g$ is positive and divisible by $b^2$. Hence $g=bg_0$, and
$b^2\mid X+1$. Use the constructive bridge, retaining its exact
$C=u_{2X+1}(A)$ and $K=u_{X+1}(P)$.

An even-index Lucas polynomial is divisible by its parameter. As $X+1$
is even and $4b^2\mid P$, we have $b^2\mid K$. In
\[
 (P-2)k=K-X-1
\]
the right side is divisible by $b^2$, while $v_2(P-2)=1$. Because $b$ is
a power of two, $b^2/2\mid k$, hence $b\mid k$. Reduction modulo three,
using $3\mid K,P,X$, gives $k\equiv2\pmod3$. Since $b\equiv1\pmod3$,
$k/b\equiv2\pmod3$, so $k=b(3k_0+2)$ for an integer $k_0$.

The recurrence also gives
$u_{2j+1}(A)\equiv(-1)^j\pmod{A^2}$. As $X$ is odd,
$A^2\mid C+1$, while $2b^2\mid B_*+1$. Thus $2b^2\mid C-B_*$.
The constructive $A$ is a multiple of four, so $v_2(A-2)=1$ and
$h=(C-B_*)/(A-2)$ is divisible by $b^2$, hence by $b$.
The normalizations in~\eqref{eq:normalization} therefore retain the
actual constructive witnesses; we have not postulated arbitrary
cofinality of Pell witnesses in chosen residue classes.

We have $\Delta>0$, $S\mid T$, and $b,Q$ positive squares.
Lemma~\ref{lem:guard} supplies a positive $z$ making $a$ a square.
Invoke the interpolation theorem to choose $\xi$ divisible by this $z$.
Write $q=T/S$, choose a root $c$ of $\pi$ modulo $\xi$, and request
$y\equiv c-q\pmod\xi$. The theorem supplies such a $y$ with $d$ a
square. Put $x=\xi/z$. Now $\Sigma\mid\Theta$, and the integral choice
\[
 m=\Theta/\Sigma+\sqrt a+\sqrt d
\]
makes~\eqref{eq:ten-final} vanish. This completes both directions.

\subsection{Exact degree}

All ten supplied variables have degree one; every fixed port has degree
zero. The decisive register degrees are
\begin{center}\small
\begin{tabular}{@{}lrrrrrr@{}}
\toprule
Register &$X$&$Y$&$A$&$C$&$K$&$Q$\\
Degree &22,964&11,863&46,694&46,696&81,524&326,092\\
\midrule
Register &$\Delta$&$a$&$d$&$\Sigma$&$Z$&$F_{10}$\\
Degree &187,881&840,069&4,109,198&233,472&233,490&\TenDegree\\
\bottomrule
\end{tabular}
\end{center}
In the inner norm expression $\Sigma^2d$ uniquely has largest degree,
4,576,142. In the outer expression $\Sigma^4d^2$ dominates the
subtracted $4\Sigma^4ad$, whose degree is only 5,883,155. This proves the
claimed exact degree once the leading form is nonzero. The independent
checker uses the actual fixed payloads and computes a nonzero coefficient
at the upper degree on the full diagonal, modulo three specified primes.
Already one such nonzero residue proves that the integral leading
coefficient is not zero; three are retained as independent transcription
controls. This is an exact lower-bound certificate, not a random
evaluation heuristic.

\section{Nine positive variables: congruence-separated norms}\label{sec:nine}

\subsection{Source and Jones registers}

For the positive-coordinate form the supplied variables are
\[
 (\epsilon_+,g_+,h_+,s,w,\phi_+,i,j,n_+)\in\Pos^9.
\]
Set $b=\epsilon_++1$, $g=2g_+$, $h=2h_+$,
$\phi=\phi_+-1$, $n=n_+-1$, and $B=a_cb^2$. Use
the common registers~\eqref{eq:source-registers}, and put
\begin{align}
 L_c&=2bB^{5741},&
 R&=L_c^2+(S_c+T_c)L_c+L_c-1-S_c,\notag\\
 N&=2^{16012040}B^{2966},& G_c&=32bB^{183}.
 \label{eq:nine-source}
\end{align}
The calibrated block length $L_c$ and the ratio threshold $N$ are
different. When $b=2^s$, $B=2^v$ with $v=\kappa+2s$, the population
calculation gives threshold exponent
\[
 5947v-30s-20=\log_2(N^2).
\]
This identity is the reason a polynomial square threshold is available
without another variable. The source estimates prove $R,N\ge8$ on all
positive assignments, even before the margin imposes an upper bound on
$g$. In fact $J_c,S_c,T_c$ are positive and $R>L_c^2$.

The retained Jones ratio registers are
\begin{align}
 U&=N^2w,&Y&=N^2s,& M_0&=RY,& P&=2M_0^2U,\notag\\
 K&=R+1+h(P-1),& A&=M_0(U+1),&B_1&=2R+1,&W_1&=bw,\notag\\
 C&=B_1+W_1+\phi,&J&=4A-5,&Q&=3W_1C-2(W_1^2-1),\notag\\
 D&=(A^2-1)C^2+1,&E&=iJDC^2,&F&=(A^2-1)E^2+1,\notag\\
 G&=A+F(CD+1-A),&H&=B_1+jC,&I&=(G^2-1)H^2+1.
 \label{eq:nine-ratio}
\end{align}
Define
\begin{align}
 a_0&=(P^2-1)K^2+1,& u&=JF,& v&=QF+J(H-C),\notag\\
 \mu&=G_c\{K^2-4(C-KY)^2\}-8gK^2,&
 M_d&=Q+H-C.\label{eq:nine-margin}
\end{align}
The imported ratio/source theorem, including its constructive direction,
equates source existence with the four conditions
\begin{equation}\label{eq:four-ratio}
 a_0\in\Sq,\quad DFI\in\Sq,\quad \mu>0,\quad u\mid v
\end{equation}
on these coordinates. The source change in Section~\ref{sec:source} meets
its quantitative hypotheses. For instance $\mu>0$ gives
$g<G_c/8=4bB^{183}$, exactly the bound used there. The retained
\code{Ratio.lean} and \code{NineRatio.lean}, and their accompanying
ordinary source proofs, specify this interface. This branch does not
depend on Sun's 2026 ten-integer theorem.

\subsection{Why the even restrictions preserve completeness}

On every supplied assignment, $B,N_0,L_c,g,S_c$ are even, so $R$ is
odd. The integer $P$ is even and $h$ is even, hence $K$ is even. Thus
\begin{equation}\label{eq:nine-parity}
 a_0\equiv1\pmod4.
\end{equation}
The fixed square multiplier used below makes the other radicand zero
modulo four. This gives noncollision without an extra multiplier involving
$h$ or a large separation padding term.

For completeness, canonical source words are even because every $w_i$ is
positive. The constructive ratio proof chooses
$K=\psi_P(R+1)$, where
$\psi_P(0)=0,\psi_P(1)=1$ and
$\psi_P(n+2)=2P\psi_P(n+1)-\psi_P(n)$. Reduction modulo two gives
$\psi_P(n)\equiv n\pmod2$. As $R+1$ is even, $K$ is even. In
$K=R+1+h(P-1)$ the coefficient $P-1$ is odd, so its positive $h$ is
already even. The constructive witnesses are therefore represented by
positive $g_+,h_+$, as required.

\subsection{A common finalizer with strict polynomial majorants}

All Jones registers used here are positive; in particular $A\ge33280$,
$u,v>0$, $H-C>0$, and $0<v/u\le M_d$. Let $\lambda$ be either 8 or 2048,
and choose polynomial integer majorants satisfying
\[
 V_0=PK>\sqrt{a_0},\qquad V_1>\lambda\sqrt{DFI}.
\]
Set
\begin{align}
 a_1&=\lambda^2DFI,& \Omega&=u(2\mu-1),\notag\\
 E_0&=un+v-\Omega M_d,&
 \Theta&=E_0-\Omega(V_0+V_1).
 \label{eq:nine-final-registers}
\end{align}
The equation is
\begin{equation}\boxed{
 F_9=\left(\Theta^2-\Omega^2(a_0+a_1)\right)^2
       -4\Omega^4a_0a_1=0.
 }\label{eq:nine-final}
\end{equation}

\begin{proposition}\label{prop:nine-final}
On the supplied positive domain, existence of $n_+$ satisfying
\eqref{eq:nine-final} is equivalent to the four conditions
\eqref{eq:four-ratio}.
\end{proposition}
\begin{proof}
The integer $2\mu-1$ is nonzero, so $\Omega\ne0$ unconditionally.
Equations~\eqref{eq:nine-parity} and $a_1\equiv0\pmod4$ give
$a_0\ne a_1$. Lemma~\ref{lem:norm} produces signed integer roots
$r_0,r_1$ with $\Theta/\Omega=r_0+r_1$. A square $a_1=\lambda^2DFI$
also makes $DFI$ a square by the rational-square-root lemma.
Since the integer majorants are strict, $|r_i|\le V_i-1$. Therefore
\[
 Z=M_d+V_0+V_1+r_0+r_1\ge M_d+2>0.
\]
The selected factor equation becomes
\[
 n+v/u=(2\mu-1)Z.
\]
Its left side is positive; hence $\mu>0$. Its right side is integral;
hence $u\mid v$. Conversely, under~\eqref{eq:four-ratio}, take any
signed roots and the same $Z$, and set
$n=(2\mu-1)Z-v/u$. This is an integer at least two because
$v/u\le M_d$. It is an allowed natural value, and $n_+=n+1$ is positive.
Substitution makes one of the norm's factors zero.
\end{proof}

\subsection{The cheaper variant}

Put $e=iJC$. Since $D-1=(A^2-1)C^2$, one has
$F=e^2D^2(D-1)+1$. Set
\[
 \lambda=8,\qquad V_1=e(8D^2-4D-1)GH.
\]
The exact identity
\begin{equation}\label{eq:fast-defect}
 \{e(8D^2-4D-1)\}^2-64DF=e^2(8D+1)-64D>0
\end{equation}
holds because $e\ge7$ and $D\ge9$. Since $I\le G^2H^2$, it gives the
strict majorant required by Proposition~\ref{prop:nine-final}.
Its additional exact identity
\[
 V_1^2-64DFI=\{e^2(8D+1)-64D\}G^2H^2+64DF(H^2-1)
\]
exposes the leading cancellation and supplies a division-free degree
certificate. The resulting polynomial $F_{9,\mathrm f}$ has exact degree
\FastDegree\ and \FastOps\ positive-domain operations.

\subsection{The lower-degree majorant}

The following polynomial majorant is inherited from the incoming
support-majorants package~\cite{majorant}. Our change is to remove the
separation padding and its multiplier, using the unconditional congruence
separation above. Write
\begin{align*}
 e&=iJC,& L_m&=CD+1-A,& N_*&=2G+L_m,\\
 N_1&=2e^2D(D-1)L_m+3C,&N_2&=2e^2(D-1)L_m,\\
 N_3&=2e^2(L_m-C),&N_4&=2e^2C,\\
 c_m&=256i(e^2+2)+2048.
\end{align*}
Define
\begin{align}
 T_m={}&e\{(1024D^2-512D-128)N_*
            -64N_1-40N_2-28N_3-21N_4\}+c_m,\notag\\
 \lambda={}&2048,\qquad V_1=T_mH.\label{eq:majorant}
\end{align}
All four $N_j$ are the displayed polynomials; no quotient computes them.
For completeness the uniform inequality, rather than only an asymptotic
expansion, is recalled here.

Put
\begin{align*}
 t_n&=1024D^6-512D^5-128D^4-64D^3-40D^2-28D-21,\\
 d_n&=33792D^5+10560D^4+4928D^3+2464D^2+1176D+441.
\end{align*}
Direct expansion gives
$t_n^2-1024^2D^{11}(D-1)=d_n>0$, and $t_n>0$ for $D\ge1$.
The exact remainders $N_*=D^jN_j+r_j$ are
\begin{align*}
 r_1&=3-A,&r_2&=3CD+3-A,\\
 r_3&=2e^2(A-1)D^2+r_2,&
 r_4&=2e^2(1-A-C)D^3+2e^2(A-1)D^2+r_2.
\end{align*}
Thus the uncorrected polynomial part of $T_m$ equals
\[
 e\frac{t_n}{D^4}N_*+
 e\left(\frac{64r_1}{D}+\frac{40r_2}{D^2}
        +\frac{28r_3}{D^3}+\frac{21r_4}{D^4}\right).
\]
For $A\ge3$ and $C,i\ge1$,
\[
 \frac{eA}{D}<\frac92i,\quad
 \frac{eC}{D}<\frac32i,\quad
 \frac{e(A+C-1)}D<6i.
\]
The possible negative remainder is less than
$252ie^2+(765/2)i$, whereas $c_m$ exceeds this plus 2048.
Writing $X_m=e^2D^3(D-1)$, the identities
$DF=X_m+D$ and $G=(X_m/D)L_m+CD+1$, together with
$\sqrt{X_m+D}<\sqrt{X_m}+D/(2\sqrt{X_m})$, give
\[
 2048\sqrt{DF}\,G
 <1024eD^2\sqrt{1-1/D}\,N_*+2048v_m,
 \qquad v_m=\frac{CD+1}{2eD\sqrt{1-1/D}}<1.
\]
Combining the estimates proves $T_m>2048\sqrt{DF}\,G>0$ uniformly.
As $I\le G^2H^2$, this proves $V_1>2048\sqrt{DFI}$, including all
assignments with arbitrarily large witness variables.

\subsection{The second cancellation and exact degree}

The decisive degrees in the calibrated nine-variable construction are
\begin{center}\small
\begin{tabular}{@{}lrrrrrr@{}}
\toprule
Register &$D$&$e$&$G$&$H$&$P$&$K$\\
Degree &139,324&69,663&719,588&22,967&81,527&81,528\\
\midrule
Register &$V_0$&$V_1$&$\Omega$&$E_0$&$a_0$&$a_1$\\
Degree &163,055&1,090,866&791,147&814,115&326,110&2,181,732\\
\bottomrule
\end{tabular}
\end{center}
Let $\delta_0=V_0^2-a_0=K^2-1$ and
$\delta_1=V_1^2-a_1$. The majorant's exact defect identity gives
\[
 \deg\delta_0=163056,\qquad\deg\delta_1=1253160.
\]
Set $B_0=E_0^2-2E_0\Omega(V_0+V_1)+\Omega^2(\delta_0+\delta_1)$.
The two-root norm admits the exact second-cancellation identity
\begin{equation}\label{eq:second-cancel}
 F_9=B_0^2+4\Omega^2V_0V_1B_0
       +4\Omega^4(\delta_0V_1^2+a_0\delta_1).
\end{equation}
The three displayed terms have degrees, respectively,
\[
 5,670,908,\qquad5,671,669,\qquad5,509,376.
\]
The middle term is uniquely highest with a positive, nonzero diagonal
leading coefficient. A simple degree walk through the primary norm would
miss this cancellation and give an overestimate.

The supplied division-free certificate constructs $D^8\delta_1$ using
the exact majorant identity, then constructs $D^8B_0$ and $D^{16}F_9$.
It is a different polynomial from the primary equation: its exact degree
is 7,900,853. Subtracting $16\deg D=2,229,184$ proves
$\deg F_9=\NineDegree$. On the diagonal of the nine input variables,
with $a_c=2^{2134940}$ and $c_c=2^{16012040}$, the leading coefficients are
\[
 \lc(F_9)=2^{507}a_c^{2835574}c_c^{348},\qquad
 \lc(D^{16}F_9)=2^{667}a_c^{3950070}c_c^{476}.
\]
Their roles are separately labeled in the exported metadata. Cancellation
by $D^{16}$ occurs only in this proof of degree; neither circuit divides.

\section{A degree-52 equation in fourteen positive variables}\label{sec:fourteen}

This secondary result combines the incoming quotient absorber, a positive
auxiliary gap, and a stronger canonical auxiliary congruence. It is useful
when degree matters more than reaching the smallest available number of
variables. Its supplied coordinates are
\[
 (J,F,\alpha,z,f,i,T,s,w,\tau,\eta,\zeta,e,\gamma)\in\Pos^{14}.
\]
Use the same four native fixed ports and define
\begin{align*}
 q&=\beta J+1,& C&=q-F-\alpha,&X&=wq,&Y&=sq^3,\\
 E&=XY,&k&=\eta+\zeta,&c&=kY+\eta,&a&=E+Y,\\
 H&=4a+3,&\Delta&=a^2+H,&\ell&=EkY,&b_0&=X+ac,\\
 G&=1+\Delta c^2,&U&=4q^4,
\end{align*}
and
\begin{align*}
 R&=\{(q-1)(q-F)+\alpha\}(q^2-1)+(M_C+qM_F)J,\\
 r&=k-R-1,\qquad D_*=Ub_0+H(U\gamma-r),\qquad M_*=D_*^2-U^2G,\\
 I&=ic^2,\qquad A_{\rm aux}=\Delta(f^2-1),\qquad W=I^2-A_{\rm aux},\\
 V&=f^2(cT-R)-c,\qquad y=V+e,\\
 N_a&=A_{\rm aux}(V^2-y^2)+y^2,\\
 N_0&=\tau^2-\ell(\ell+k),\qquad
 N_t=(K_0+w)C+q-F-z(q-1).
\end{align*}
The equation is
\begin{equation}\boxed{
 P_{14}=M_*^2+W^2+N_a^2+(N_0N_t-1)^2-1=0.
 }\label{eq:fourteen}
\end{equation}

\subsection{A minimum block and soundness}

For $f=1$, $A_{\rm aux}=0$ and $W=I^2>1$. For $f\ge2$,
$A_{\rm aux}>1$ and $N_a$ cannot vanish. Otherwise
$A_{\rm aux}V^2=(A_{\rm aux}-1)(V+e)^2$ would make
$A_{\rm aux}(A_{\rm aux}-1)$ a rational square. Neither factor involving
$V$ can be zero because $e>0$, and the consecutive product lies strictly
between $(A_{\rm aux}-1)^2$ and $A_{\rm aux}^2$. Hence always
$W^2+N_a^2\ge1$. Equality forces $W=0$ and $N_a=\pm1$. With
$A_{\rm aux}=I^2$, reducing
$N_a=I^2V^2-(I^2-1)y^2$ modulo four excludes the negative sign.
Thus equation~\eqref{eq:fourteen} is equivalent tuplewise to
\[
 M_*=0,\quad W=0,\quad N_a=1,\quad N_0N_t=1.
\]
Furthermore $N_0=-1$ is impossible. Here $\ell=nk$ with
$n=XY^2\ge2$, so it would give
$\tau^2-n(n+1)k^2=-1$. For a solution of minimum positive $k$, multiplication
by the unit $2n+1-2\sqrt{n(n+1)}$ gives a new second coordinate
$k'=(2n+1)k-2\tau$ with $0<k'<k$, a contradiction. The inequalities
follow from $nk<\tau<(n+1/2)k$. Thus $N_0=N_t=1$.

Map a new zero to the incoming fourteen-variable construction by
$T_{\rm old}=fT$ and $y_{\rm old}=|V+e|$. The identity
\[
 V=c(T_{\rm old}f-1)-Rf^2
\]
is exact, and $y_{\rm old}>0$ follows from $N_a=1$ and $I^2>1$.
All old residuals vanish. The retained absorber proof then recovers its
omitted positive quotient before using the strong Pell rank theorem;
the order of that recovery is essential. For reference, it first derives
$C>0$, $q\ge16$, and $3q<R<q^4-2$ from transport. The modulus
$U/\gcd(U,H)$ exceeds $R+1$, because $\gcd(U,H)\le3$. First-norm
classification supplies an early index bound; only then is the strong
rank theorem applied to recover the original compiler witness.
The complete inherited proof is retained in
\code{fourteen/upstream/FOURTEEN_PROOF.md}.

\subsection{The stronger auxiliary congruence}

It remains to show that requiring $T_{\rm old}$ to be a multiple of $f$
does not lose a canonical witness. At a normalized compiler point the
actual packed index is $R\equiv3\pmod4$. Put $A=a+2$, so
$\Delta=A^2-1$, and use the ordinary Pell sequences $\chi_A,\psi_A$.
The established completion chooses
\[
 c=\psi_A(R),\quad m_0=2cR,\quad f=\chi_A(m_0),\quad
 I_0=\psi_A(m_0),\quad \rho=\Delta I_0,
\]
and $V=\chi_\rho(R)/\rho$, $y=\psi_\rho(R)$. The strong coordinate is
chosen so that $I=ic^2=\rho$. In particular
$\rho^2=\Delta(f^2-1)$ and $c^2\mid\rho$.

For $R=2r_0+1$, the odd quotient polynomial $Q_{r_0}$ satisfies
\[
 \chi_z(R)=zQ_{r_0}(z^2),\quad
 Q_{r_0}(1-A^2)=(-1)^{r_0}\psi_A(R),\quad
 Q_{r_0}(0)=(-1)^{r_0}R.
\]
Since $r_0$ is odd,
\begin{equation}\label{eq:squared-congruence}
 V=Q_{r_0}(\rho^2)\equiv-c\pmod{f^2}.
\end{equation}
Also $V\equiv-R\pmod c$ and $f^2\equiv1\pmod c$. Therefore
$o=(V+c)/f^2$ is positive integral, $c\mid o+R$, and
$T=(o+R)/c$ is positive integral. This is precisely the new formula
$V=f^2(cT-R)-c$. The norm identity gives
\[
 (\rho^2-1)(y^2-V^2)=V^2-1>0,
\]
because $R\ge3$, $\rho>1$, and $V\ge4\rho^2-3>1$.
Hence $e=y-V>0$. The retained quotient absorber then supplies its positive
carrier $\gamma$. This constructs all fourteen positive coordinates.

\subsection{Exact degree 52}

The main residual has exact degree 26. Write $D_0=b_0+H\gamma$. Then
\begin{align*}
 M_*&=U^2(D_0^2-G)-2UHrD_0+H^2r^2,\\
 D_0^2-G&=2ac(X+H\gamma)+(X+H\gamma)^2-Hc^2-1.
\end{align*}
Its leading form is $128w^2s^3(\eta+\zeta)\gamma(\beta J)^{19}$.
The strong residual $W$ has degree 22; the gap identity
\[
 N_a=V^2-(A_{\rm aux}-1)e(2V+e)
\]
gives degree at most 23, and $N_0N_t-1$ has degree at most 24.
Only $M_*^2$ reaches degree 52. Thus
\[
 (P_{14})_{52}=16384w^4s^6(\eta+\zeta)^2\gamma^2(\beta J)^{38}\ne0.
\]
The circuit has exactly 79 live operations. The stronger congruence
\eqref{eq:squared-congruence} is already implicit in the canonical Pell
completion; its use here pays for the extra gap coordinate without
increasing the fourteen-variable evaluation cost.

\section{Verification, dependencies, and formal status}\label{sec:status}

\subsection{What the archive checks}

The top-level \code{verify.py} runs the finite checks offline, using Python's
standard library. It rebuilds the literal circuits, checks their live-gate
counts, and validates the retained fixed payloads. Individual reports give
the exact check counts and hashes; the principal evidence is:
\begin{itemize}
\item Complete support isolation: all 49,632 interactions, 94 intended hits,
  and the five-collision legacy negative control. A second checker
  independently parses the residual expressions and reconstructs the
  emitted sparse polynomial, rather than trusting the builder's source
  grouping.
\item Fixed-coefficient modular evaluation of the displayed ten-variable
  registers and their universal congruences, together with exact symbolic
  identities for the two-root extraction and complete degree metadata.
\item Independent finite ring expansion of the majorant identities and
  the second-cancellation identity, plus actual-coefficient comparisons
  of the majorant certificate with $D^{16}$ times the primary equation.
\item Exact degree certificates using leading forms and nonzero
  coefficients at known upper degrees; these prove degree, rather than
  merely guessing it from a few evaluations.
\item The fourteen-variable main-residual expansion, gap identity,
  coefficient binding, and full univariate degree replay modulo three
  primes, with targeted canonical Pell and minimum-block controls.
\item Regeneration of the fixed native coefficient tuple from finite
  machine/tile data, with the inherited finite geometry and carry checks
  where supplied by its generator.
\end{itemize}

The finite checks establish finite identities, transcriptions, and
certificate claims. They do not establish unsolvability by bounded
search, and no sufficiently long simulation could do so. The unbounded
existence equivalences are the mathematical proofs in this report and the
explicitly retained inherited theorem chain.

\subsection{Three different proof levels}

\begin{enumerate}
\item \textbf{New ordinary arithmetic proofs.} The positive-word source
  transformation, the two-radical separation, the coprime guard, preservation
  of the new congruence restrictions, the nine-variable finalizer, and the
  fourteen-variable witness reconstruction are proved here. These use
  elementary integer arithmetic, classical Pell/Lucas theorems, and the
  stated source interfaces, all formalizable in ZFC.
\item \textbf{Inherited unbounded theorems.} The native coefficient compiler
  and K215 calibration identify the literal source with syntactic ZFC
  refutations. The nine-variable branch uses the Jones ratio theorem; the
  ten-integer branch additionally uses Sun's modified Lucas interpolation
  theorem from the 2026 preprint. Their use is neither a new proof of those
  results nor a claim of independent external peer review.
\item \textbf{Proof-assistant artifacts.} The pinned repository has a
  separate external Lean theorem for a closed generic nine-variable
  coefficient recipe. That theorem is not a kernel check of these new
  small literal circuits. Its status document also distinguishes the
  external equivalence from an exported object-language ZFC derivation of
  the two quoted sentences. This release does not supply a new end-to-end
  Lean/Rocq certificate or such an object-language derivation.
\end{enumerate}

The native bundle retains the relevant exact source files, licenses and
source pins, including the portable K215 endpoint source closure where
available. Rebuilding that Lean endpoint requires its stated toolchain;
it was not rerun in this session. The separate source-chain ledger records
this boundary. The independent reviews in this archive are additional
AI-assisted reviews within the session, not external human peer review.

\subsection{Reproduction and integration}

From an extracted archive, run
\begin{verbatim}
python3 verify.py
pdflatex -interaction=nonstopmode -halt-on-error report.tex
pdflatex -interaction=nonstopmode -halt-on-error report.tex
\end{verbatim}
The first command is the finite arithmetic audit; TeX is needed only to
rebuild the report. The exact per-component commands and output names are
listed in \code{README.md}. The manifest records hashes for delivered
files. Generated certificate circuits are clearly distinguished from the
recommended primary circuits. No network access or search for a refutation
is part of the polynomial generators.

For repository incorporation, retain this as a new incoming construction
and update the frontier by domain and metric. Its natural first integration
targets are the source extraction lemma, the universal residue identities,
the two-root norm lemma, and the guard theorem. These are compact formal
proof obligations. An end-to-end formal claim should be made only after
connecting the literal emitted coefficient maps to the retained semantic
theorems and exporting the desired object-language formulation if required.

\section{What is new, and what remains open}

The main improvement is a structural reduction in the number of radicands
that must be combined. The ten-integer equation no longer carries the base
square, an auxiliary Lucas square, and positivity as separate radical
requirements. The coprime product $r=bQ$, together with residues modulo
three, recovers both square conditions from one square guard. Residues
modulo the even base distinguish this guard from the Lucas certificate.
A norm with four factors therefore replaces the earlier eight-factor norm.
This is responsible for most of the ten-variable degree reduction.

The nine-positive reduction exploits the same separation principle with
an easier parity normalization. It can discard both the extra multiplicative
separator and the large gap padding of the old majorant finalizer. The
second defect expansion then exposes a further leading cancellation. The
source transformation is complementary: it reduces the number of fixed
source records, improving circuit costs for both branches without
increasing the target span. It requires the repaired finite support
certificate; merely reusing the earlier targets would be wrong.

The fourteen-variable result is a smaller local improvement. Its stronger
auxiliary congruence and positive gap lower the degree of the residual that
previously dominated, leaving the main norm as the unique degree-26 term
before squaring. It does not compete with the few-variable constructions
on variable count.

Several next directions are precise but unproved here. A shorter target
code for the 94-record source could reduce both few-variable degrees.
Further radicand combination would need a new arithmetic identity; the
two-root lemma alone does not justify combining three radicals without
separation. Reducing a supplied variable would require preserving
constructive congruence freedom through the Lucas interpolation or finding
a different source representation. Finally, a smaller coherent native
tuple would improve coefficient height, a metric untouched by the present
word transformation. These are research targets, not claimed results.

\appendix
\section{The complete positive quadratic source}\label{app:source}

The 31 coordinates, in the order used by the weight vector, are
\begin{gather*}
 (J,F,\alpha,z,f,h,i,o,j,s,w,\tau,\eta,\zeta,y,\gamma,\\
 N,E,P,c,a,X,D,C_2,L,F_2,V,g,I,B_{\rm aux},\Delta).
\end{gather*}
Set $q=\beta J+1$, $k=\eta+\zeta$, $C=q-F-\alpha$; these are affine
abbreviations. The residual is left side minus right side in each row.
\begin{longtable}{@{}r>{\raggedright\arraybackslash}p{0.87\textwidth}@{}}
\toprule Row & Equation\\\midrule\endhead
1 & $L=q^2-1$\\
2 & $N=q(L+1)$\\
3 & $X=wq$\\
4 & $a=E+sN$\\
5 & $E=X(a-E)$\\
6 & $c=k(a-E)+\eta$\\
7 & $P=E(c-\eta)+k$\\
8 & $\tau^2=P(P-k)+1$\\
9 & $D=X+ac+\gamma(4a+3)$\\
10 & $\Delta=a^2+4a+3$\\
11 & $C_2=c^2$\\
12 & $D^2=\Delta C_2+1$\\
13 & $k=jc-V+1+hE$\\
14 & $jc-V=(q-F)(N-q)-CL+(M_C+qM_F)J$\\
15 & $(K_0+w)C+(q-F)-z(q-1)=1$\\
16 & $F_2=f^2$\\
17 & $V=of-c$\\
18 & $g=y^2-V^2$\\
19 & $I=iC_2$\\
20 & $B_{\rm aux}=I^2$\\
21 & $B_{\rm aux}=\Delta(F_2-1)$\\
22 & $y^2-B_{\rm aux}g=1$\\
\bottomrule
\end{longtable}
The original sparse JSON has SHA-256
\begin{center}\small\ttfamily
1c336ba0b3c9827ffddcf7ce895cad9dca2907f907cde7484c7f387bb2cb1287.
\end{center}
The plain expression parser and the sparse-record parser independently
reconstruct the same polynomial data. The shifted transport convention
of a different package is not substituted into row 15.

\clearpage
\section{Literal selective code and fixed coefficient recipe}\label{app:code}

The weights are
\begin{align*}
 (w_1,\ldots,w_{31})={}&(9,25,23,139,35,151,168,21,91,181,94,52,150,180,7,135,\\
 &38,88,183,106,178,136,20,138,8,64,105,154,37,137,120).
\end{align*}
The corrected targets, tabulated by original source row, are
\begin{center}\small
\begin{tabular}{@{}rrrrrrrrrrr@{}}
\toprule
Row&1&2&3&4&5&6&7&8&9&10\\
Target&2792&86&2358&3929&5190&4238&1296&3634&451&3157\\
\midrule
Row&11&12&13&14&15&16&17&18&19&20\\
Target&4384&4865&2124&685&1530&70&3290&4527&2696&2780\\
\midrule
Row&21&22&&&&&&&&\\
Target&1059&1874&&&&&&&&\\
\bottomrule
\end{tabular}
\end{center}
Together with the 22 residuals and the factorial formula
\eqref{eq:Dcode}, these tables specify every coefficient of $\mathcal D$.
The sorted target order is retained in \code{code_positive.json}.

For completeness, the native tuple's finite numerical recipe is as follows.
Let $\rho=2^{121}$, $B_n=\rho^{8801}$, and
\begin{align*}
 \mathcal E={}&\{0,\ldots,942\}\cup\{983,\ldots,992\}
 \cup\{1003,\ldots,1032\}\\
 &\cup\{1063,1064,1095,1096,2195,3296\}.
\end{align*}
The retained literal 943 tiles have four ordered slot labels $x_{it}$.
For $0\le j<44$ put $W_j=2^{10+2j}$. At tile selector $i$, set
\[
 c_i=1+\sum_{t=0}^3\sum_{j=0}^9
 \operatorname{bit}_j(x_{it}+512)W_{10t+j}
 +W_{40}+W_{41}+W_{42}+W_{43}.
\]
At copy exponents $o_t+j$, where $(o_0,o_1,o_2,o_3)=(1023,1013,1003,983)$,
put $c_{o_t+j}=W_{10t+j}$. At anchors $1095,1096,2195,3296$ put,
respectively, $W_{40},W_{41},W_{42},W_{43}$. Set $c_{1063}=c_{1064}=0$.
Define
\begin{align*}
 \mu_n&=1022+\sum_{j=0}^{43}W_j,\\
 Z_n&=(2^{119}-1)\sum_{e=8792}^{8798}\rho^e+(2^{27}-1)\rho^{8799},\\
 D_C&=\rho^{30}+\rho^{62}+\rho^{1100}+\rho^{2200}
       +\sum_{e\in\mathcal E}c_e(\rho^{4393-e}+\rho^{5495-e}),\\
 D_R&=\rho^{42},\\
 M_C&=B_n-1-\sum_{e\in\mathcal E\setminus\{1\}}\rho^e-2\rho^{1063},\\
 M_F^{\rm native}&=4+\mu_n(\rho^{4393}+\rho^{5495})
   +\sum_{j=0}^{9}(\rho^{1013+j}+\rho^{1065+j})+\rho^{2195}+\rho^{3296}+Z_n.
\end{align*}
The four ports are
\[
 \beta=B_n-1,\qquad K_0=D_C+B_nD_R,\qquad
 M_C,\qquad M_F=M_F^{\rm native}+B_n-1.
\]
The tile list is regenerated by bounded arithmetic from the 430 literal
actions of the retained K215 machine. The finite native generator and its
source-chain ledger specify the exact ordering and all byte formats.
None of these calculations requires knowing whether the machine halts.

\begin{thebibliography}{9}
\bibitem{repo}
V. Reshetnikov, \emph{ProveIt}, repository snapshot
\texttt{4b4349064a750ba6347293b55874a7f92abdd0db}, 23 October 2026.
\url{https://github.com/VladimirReshetnikov/ProveIt/tree/4b4349064a750ba6347293b55874a7f92abdd0db/Computability/DiophantineConsistency/con-zfc-diophantine-forms}.

\bibitem{selective}
\emph{ZFC Diophantine Selective Quadratic}, incoming archive at the same
repository snapshot, 23 October 2026. Baseline nine-positive and
ten-integer constructions; selective extraction, sparse borrows, and
calibrated source proofs.
\url{https://github.com/VladimirReshetnikov/ProveIt/blob/4b4349064a750ba6347293b55874a7f92abdd0db/docs/incoming/ZFC_Diophantine_Selective_Quadratic_2026-10-23.zip}.

\bibitem{majorant}
\emph{Con-ZFC support majorants}, incoming archive, 23 October 2026.
The retained \code{HIGHER_MAJORANT.md} contains the polynomial square-root
majorant and its exact defect identity.
\url{https://github.com/VladimirReshetnikov/ProveIt/blob/4b4349064a750ba6347293b55874a7f92abdd0db/docs/incoming/con-zfc-support-majorants-2026-10-23.zip}.

\bibitem{support}
\emph{ConZFC Support Sensitive}, incoming archive, 23 October 2026.
Fourteen-variable quotient absorber and its ordered index-recovery proof.
\url{https://github.com/VladimirReshetnikov/ProveIt/blob/4b4349064a750ba6347293b55874a7f92abdd0db/docs/incoming/ConZFC_Support_Sensitive_2026-10-23.zip}.

\bibitem{sun2021}
Z.-W. Sun, \emph{Further results on Hilbert's Tenth Problem},
Science China Mathematics \textbf{64} (2021), 281--306;
arXiv:1704.03504v7, especially Theorems 4.12--4.13 and equation (4.16).
\url{https://arxiv.org/html/1704.03504v7}.

\bibitem{sun2026}
Z.-W. Sun, \emph{Ten unknowns for Hilbert's tenth problem over the integers},
arXiv:2609.01594v2, 2 September 2026, especially Theorems 3.1 and 4.1 and
Propositions 6.1--6.2. Preprint dependency of the ten-integer branch.
\url{https://arxiv.org/html/2609.01594v2}.

\bibitem{jones}
J. P. Jones, \emph{Universal Diophantine equation},
Journal of Symbolic Logic \textbf{47} (1982), 549--571.
The exact ratio interface used here is also preserved in the pinned
repository's \code{Ratio.lean} and \code{NineRatio.lean} sources.
\end{thebibliography}
\end{document}
